\subsection{Transporte de la sección electrónica al sector electrodébil}
\label{amp-ew-seccion}

La sección central de \eqref{md:hib:electron} conserva conjuntamente el
prefactor, el corrector y la memoria de retorno. Sus dependencias proceden
de las hojas APP, de la orientación TRIT y de la ruta registrada por TPK;
el carácter angular y las torres actúan después sobre las multisecciones
de la misma estructura discreta del continuo. Escribimos
\begin{equation}
 r_e=R_{\rm act}^{\kappa_e},\qquad
 E_e^\beta=\mathcal U_E\varepsilon_e^{(0)}r_e,\qquad
 \kappa_e=80-54\alpha+6\Delta_4.
 \label{amp-ew-electron}
\end{equation}
Aquí $\mathcal U_E$ es la sección de la recta energética y $r_e$ abrevia
el retorno electrónico. La década de acción permanece en su recta propia;
la coordenada energética realizada se transporta una sola vez.

Las firmas angulares de \eqref{amp-higgs-firmas} producen, en la misma
coordenatización, $E_s^\angle=E_e^\beta\exp(n_sx+s_sy/6)$. En particular,
\begin{equation}
 E_W^\angle=\mathcal U_E\varepsilon_e^{(0)}r_e
                  e^{94x-y/6},\qquad D=e^{-2x}.
 \label{amp-ew-w}
\end{equation}
Se utiliza esta realización angular de W. Los refinamientos de otra torre
conservan sus factores propios cuando se compone con ellos.

\begin{proposition}[Escala electrónica y acoplamiento de Fermi de árbol]
En la realización electrodébil de árbol, con las coordenadas energéticas
expresadas en una misma unidad y $g^2=4\pi\alpha/(1-D)$, las relaciones
$v=2E_W^\angle/g$ y $G_F^{(0)}=1/(\sqrt2v^2)$ dan
\begin{equation}
 G_F^{(0)}=
 \frac{\pi\alpha\,e^{-188x+y/3}}
 {\sqrt2(1-D)\mathcal U_E^2(\varepsilon_e^{(0)})^2r_e^2}.
 \label{amp-ew-fermi}
\end{equation}
Con las firmas, los canales y la coordenatización energética fijos, el paso de la
etapa basal a la etapa retornada satisface
\begin{equation}
 \frac{(G_F^{(0)})_\beta}{(G_F^{(0)})_0}=r_e^{-2},\qquad
 \frac{v_\beta}{v_0}=r_e.
 \label{amp-ew-escala}
\end{equation}
Los cocientes $E_W^\angle/E_Z^\angle$, $E_H^\angle/E_W^\angle$ y
$\lambda_H=(g^2/8)(E_H^\angle/E_W^\angle)^2$ son independientes de
ese factor común $r_e$.
\end{proposition}
\begin{proof}
La sustitución de $v$ da $G_F^{(0)}=g^2/[4\sqrt2(E_W^\angle)^2]$.
Aplicar la expresión de $g^2$ y elevar \eqref{amp-ew-w} al cuadrado
produce \eqref{amp-ew-fermi}: los coeficientes $188$ y $1/3$ resultan
respectivamente de $2\cdot94$ y $2/6$. Cambiar $E_e^0$ por
$E_e^\beta=r_eE_e^0$ multiplica todas las energías angulares por $r_e$.
La definición de $v$ conserva la primera potencia; la de $G_F^{(0)}$
conserva la inversa de su cuadrado. En cada cociente de energías se
cancela el mismo factor positivo, y la fórmula de $\lambda_H$ conserva
esa cancelación. Los acoplamientos $g,g'$ dependen de $\alpha,x$ y
permanecen fijos en esta comparación.
\end{proof}

La magnitud $G_F^{(0)}$ tiene dimensión de energía inversa al cuadrado.
La realización de árbol declarada especifica el alcance físico de esta
composición; las correcciones radiativas tienen sus operadores y su orden
posteriores. Comparar las dos etapas del lector conserva sus coordenatizaciones y
sus firmas, en lugar de postular una variación temporal independiente de
las constantes correlacionadas.

\subsection{Conservación de los refinamientos de ruta}
\label{amp-ew-refinamiento}
La composición con la ley completa preserva el carácter y la torre que
corresponden a cada ruta. Para dos autosecciones positivas comparables,
en una misma coordenatización y con un mismo lector $r\in\{D,\Omega\}$, sea
\begin{equation}
 \rho_{W/e}^{\,r}
 =\frac{\chi_{\rm full}^{\,r}(\sigma_W,\eta_W)}
        {\chi_{\rm full}^{\,r}(\sigma_e,\eta_e)}
 R_{\rm act}^{\kappa_r(\Gamma_W)-\kappa_e}.
 \label{amp-ew-razon-completa}
\end{equation}
La energía correspondiente es $E_W=E_e^\beta\rho_{W/e}^{\,r}$.
En la realización de Fermi compatible con esa sección se obtiene
\begin{equation}
 G_F^{(0)}=
 \frac{\pi\alpha}
 {\sqrt2(1-D)(E_e^\beta)^2(\rho_{W/e}^{\,r})^2}.
 \label{amp-ew-fermi-completo}
\end{equation}
La prueba sustituye la razón completa en la identidad cuadrática
anterior. En la coordenatización angular explícita, esa razón adopta la forma de
\eqref{amp-ew-w}; en una realización refinada retiene también las
diferencias de torres y de lectores de \eqref{md:hib:cocientes}.
Si $E_s=E_e^\beta\chi_sf_s$, entonces
$E_s/E_t=(\chi_s/\chi_t)(f_s/f_t)$. Esta igualdad muestra exactamente
qué factor se cancela y qué información específica permanece.

\subsection{Congruencia másica e inversión espectral sin conmutación}
\label{amp-ew-congruencia}
Sea $\widehat M^0>0$ el operador basal sobre una fibra finita positiva,
con su carácter y sus torres incorporados, y sea $\widehat B^r$
el lector autoadjunto de retorno. El transporte de acción produce
$C=R_{\rm act}^{\widehat B^r/2}>0$ y la ley
$\widehat M^\beta=C\widehat M^0C$.

\begin{proposition}[Reciprocidad operatoria de la longitud de acción]
Con $\hbar_a,c>0$ en una misma coordenatización dimensional,
\begin{equation}
 \begin{aligned}
 (\widehat M^\beta)^{-1}
    &=C^{-1}(\widehat M^0)^{-1}C^{-1},\\
 \widehat L_C^\beta:=\frac{\hbar_a}{c}(\widehat M^\beta)^{-1}
    &=C^{-1}\widehat L_C^0C^{-1},\qquad
 \widehat L_C^0=\frac{\hbar_a}{c}(\widehat M^0)^{-1}.
 \end{aligned}
 \label{amp-ew-inversa-operador}
\end{equation}
Estas identidades admiten $[C,\widehat M^0]\ne0$.
\end{proposition}
\begin{proof}
La positividad de $C$ implica invertibilidad. Para $u\ne0$,
$\langle u,C\widehat M^0Cu\rangle
=\langle Cu,\widehat M^0Cu\rangle>0$; por tanto, el operador final
también es invertible. La multiplicación directa en ambos órdenes da
\[
 (C\widehat M^0C)\bigl(C^{-1}(\widehat M^0)^{-1}C^{-1}\bigr)=I,
 \qquad
 \bigl(C^{-1}(\widehat M^0)^{-1}C^{-1}\bigr)(C\widehat M^0C)=I.
\]
Multiplicar por $\hbar_a/c$ prueba la segunda identidad. En la
descomposición espectral del operador final,
$\widehat M^\beta=\sum_jm_jP_j$ con $m_j>0$, se obtiene
$\widehat L_C^\beta=\sum_j\hbar_a/(cm_j)P_j$. La función recíproca
es inyectiva sobre los autovalores positivos: conserva sus proyectores
y multiplicidades finales.
\end{proof}

La congruencia puede cambiar los proyectores del operador basal; la
inversión espectral conserva los del operador final. El sector nulo
queda fuera de esta inversión. Una preparación que ocupe varios
autoespacios conserva su medida espectral y sus multiplicidades.
En rango uno, la identidad recupera la longitud electrónica
$\bar\lambda_{C,e}=\hbar_ac/E_e^\beta$ y, después,
$a_{{\rm B},e}=\bar\lambda_{C,e}/\alpha$. El resultado conecta la
escala electrónica, la respuesta electrodébil y la longitud de acción
mediante operaciones explícitas sobre sus estructuras compartidas.
